%=====================================================================
%  preamble.tex — Cấu hình chuẩn cho Báo cáo Đồ án Tốt nghiệp Bách Khoa
%  Biên dịch: pdflatex (chạy 2-3 lần để cập nhật mục lục / tham chiếu)
%=====================================================================

%---------------------- Ngôn ngữ & phông chữ -------------------------
\usepackage[utf8]{vntex}
\usepackage{mathptmx}          % Times New Roman cho cả văn bản và công thức

%---------------------- Trình bày trang ------------------------------
\usepackage{geometry}
\geometry{a4paper, left=30mm, right=20mm, top=25mm, bottom=25mm}
\usepackage{setspace}
\onehalfspacing                % Giãn dòng 1.5 theo quy định báo cáo
\setlength{\parindent}{1cm}
\usepackage{indentfirst}
\usepackage{lastpage}
\usepackage{emptypage}

%---------------------- Bảng biểu & hình ảnh -------------------------
\usepackage{graphicx}
\graphicspath{{Images/}{Images/Chuong2/}{Images/Chuong3/}{Images/Chuong4/}{Images/Chuong5/}{Images/Chuong6/}{Images/GiaoDien/}}
\usepackage{array, multirow, multicol, longtable, tabularx, booktabs, hhline}
\usepackage{float}
\usepackage{caption, subcaption}
\usepackage{rotating}
\usepackage{diagbox}
\usepackage{tikz}
\usetikzlibrary{arrows, arrows.meta, backgrounds, calc, positioning, shapes.geometric}

%---------------------- Toán học & thuật toán ------------------------
\usepackage{amsmath, amssymb, amsthm}
\usepackage[lined, boxed, commentsnumbered]{algorithm2e}

%---------------------- Mã nguồn -------------------------------------
\usepackage{xcolor}
\usepackage{listings}
\lstdefinestyle{cospace}{
    basicstyle=\footnotesize\ttfamily,
    keywordstyle=\color{blue!70!black}\bfseries,
    commentstyle=\color{green!40!black}\itshape,
    stringstyle=\color{red!60!black},
    numbers=left, numberstyle=\tiny\color{gray}, numbersep=8pt,
    frame=single, rulecolor=\color{gray!50},
    backgroundcolor=\color{gray!5},
    breaklines=true, breakatwhitespace=false,
    showstringspaces=false, tabsize=2, captionpos=b,
    extendedchars=true, literate={ă}{{\u a}}1 {â}{{\^a}}1 {đ}{{\dj}}1 {ê}{{\^e}}1 {ô}{{\^o}}1 {ơ}{{\ohorn}}1 {ư}{{\uhorn}}1
}
\lstset{style=cospace}

%---------------------- Danh mục & liên kết --------------------------
\usepackage[unicode]{hyperref}
\hypersetup{
    colorlinks=true,
    linkcolor=black, citecolor=black, urlcolor=blue,
    pdftitle={Bao cao Do an Tot nghiep - He thong quan ly Co-working Space CoSpace},
    pdfauthor={Vuong Song Toan - Le Dang Khoa}
}
\usepackage[nottoc]{tocbibind}   % đưa Tài liệu tham khảo vào mục lục

%---------------------- Việt hóa tên mục ----------------------------
\renewcommand{\chaptername}{Chương}
\renewcommand{\contentsname}{Mục lục}
\renewcommand{\listfigurename}{Danh sách hình vẽ}
\renewcommand{\listtablename}{Danh sách bảng}
\renewcommand{\bibname}{\texorpdfstring{Tài liệu tham khảo}{Tai lieu tham khao}}
\renewcommand{\appendixname}{Phụ lục}
\renewcommand{\figurename}{Hình}
\renewcommand{\tablename}{Bảng}
\renewcommand{\lstlistingname}{Đoạn mã}

%---------------------- Đầu trang / chân trang -----------------------
\usepackage{fancyhdr}
\setlength{\headheight}{40pt}
\pagestyle{fancy}
\fancyhf{}
\fancyhead[L]{%
 \begin{tabular}{rl}
    \begin{picture}(25,15)(0,0)
    \put(0,-8){\includegraphics[width=10mm, height=10mm]{hcmut.png}}
   \end{picture}&
	\begin{tabular}{l}
		\textbf{\small Trường Đại học Bách khoa -- Đại học Quốc gia TP.HCM}\\
		\textbf{\small Khoa Khoa học \& Kỹ thuật Máy tính}
	\end{tabular}
 \end{tabular}%
}
\fancyfoot[L]{\scriptsize Báo cáo Đồ án Tốt nghiệp -- Hệ thống quản lý văn phòng chia sẻ CoSpace}
\fancyfoot[R]{\scriptsize Trang {\thepage}}
\renewcommand{\headrulewidth}{0.3pt}
\renewcommand{\footrulewidth}{0.3pt}
\fancypagestyle{plain}{%   % trang đầu mỗi chương vẫn giữ header/footer
  \fancyhf{}
  \fancyhead[L]{%
    \begin{tabular}{rl}
      \begin{picture}(25,15)(0,0)\put(0,-8){\includegraphics[width=10mm, height=10mm]{hcmut.png}}\end{picture}&
      \begin{tabular}{l}
        \textbf{\small Trường Đại học Bách khoa -- Đại học Quốc gia TP.HCM}\\
        \textbf{\small Khoa Khoa học \& Kỹ thuật Máy tính}
      \end{tabular}
    \end{tabular}%
  }
  \fancyfoot[L]{\scriptsize Báo cáo Đồ án Tốt nghiệp -- Hệ thống quản lý văn phòng chia sẻ CoSpace}
  \fancyfoot[R]{\scriptsize Trang {\thepage}}
}

% Cấu hình footer riêng cho phần nội dung chính (Đánh số Ả Rập: Trang x/y)
\newcommand{\setupArabicFooters}{%
  \fancyfoot[R]{\scriptsize Trang {\thepage}/\pageref{LastPage}}%
  \fancypagestyle{plain}{%
    \fancyhf{}%
    \fancyhead[L]{%
      \begin{tabular}{rl}%
        \begin{picture}(25,15)(0,0)\put(0,-8){\includegraphics[width=10mm, height=10mm]{hcmut.png}}\end{picture}&%
        \begin{tabular}{l}%
          \textbf{\small Trường Đại học Bách khoa -- Đại học Quốc gia TP.HCM}\\%
          \textbf{\small Khoa Khoa học \& Kỹ thuật Máy tính}%
        \end{tabular}%
      \end{tabular}%
    }%
    \fancyfoot[L]{\scriptsize Báo cáo Đồ án Tốt nghiệp -- Hệ thống quản lý văn phòng chia sẻ CoSpace}%
    \fancyfoot[R]{\scriptsize Trang {\thepage}/\pageref{LastPage}}%
  }%
}

%---------------------- Cấp mục & đánh số ---------------------------
\setcounter{secnumdepth}{4}
\setcounter{tocdepth}{3}
\makeatletter
\newcounter{subsubsubsection}[subsubsection]
\renewcommand\thesubsubsubsection{\thesubsubsection.\@alph\c@subsubsubsection}
\newcommand\subsubsubsection{\@startsection{subsubsubsection}{4}{\z@}%
    {-3.25ex\@plus -1ex \@minus -.2ex}{1.5ex \@plus .2ex}{\normalfont\normalsize\bfseries}}
\newcommand*\l@subsubsubsection{\@dottedtocline{4}{10.0em}{4.1em}}
\providecommand*{\toclevel@subsubsubsection}{4}
\newcommand*{\subsubsubsectionmark}[1]{}
\makeatother

%---------------------- Định lý / định nghĩa -------------------------
\newtheorem{definition}{\bf Định nghĩa}[chapter]
\newtheorem{theorem}{\bf Định lý}[chapter]
\newtheorem{property}{\bf Tính chất}[chapter]

%=====================================================================
%  CÁC LỆNH TIỆN ÍCH DÙNG TRONG BÁO CÁO
%=====================================================================

% Bật/tắt khung ghi chú (mặc định tắt khi hoàn thiện)
\newif\ifhienghichu
\hienghichufalse

\newcommand{\ghichu}[1]{%
  \ifhienghichu
    \par\vspace{2pt}
    \noindent\begin{tikzpicture}
      \node[draw=gray!55, fill=gray!8, rounded corners=3pt, line width=0.4pt,
            inner sep=7pt, text width=\dimexpr\linewidth-18pt\relax, align=justify]
        {\small\itshape\color{black!72}\textbf{[Ghi chú]}~#1};
    \end{tikzpicture}
    \par\vspace{4pt}
  \fi
}

\newcommand{\mtund}{\textunderscore\allowbreak}
\newcommand{\mt}[1]{{\let\_\mtund\texttt{#1}}}

\makeatletter
\def\nguon@bs#1/#2\nguon@stop{%
  #1%
  \ifx\relax#2\relax\else /\allowbreak\nguon@bs#2\nguon@stop\fi}
\newcommand{\nguon@brk}[1]{\nguon@bs#1/\nguon@stop}

\newcommand{\nguon}[1]{%
  \ifhienghichu
    \par\noindent{\raggedright\footnotesize\color{blue!45!black}\textbf{Nguồn dữ liệu:} \texttt{\let\_\mtund\nguon@brk{#1}}\par}\vspace{4pt}
  \fi
}
\makeatother

% Chèn sơ đồ UML thẳng đứng
\newcommand{\hinhuml}[3]{%
  \begin{figure}[H]
    \centering
    \includegraphics[width=\linewidth,height=0.72\textheight,keepaspectratio]{#1.png}
    \caption{#2}
    \label{#3}
  \end{figure}
}

% Chèn sơ đồ UML hiển thị xuôi chiều trang (phù hợp đọc trên màn hình máy tính)
\newcommand{\hinhumlngang}[3]{%
  \begin{figure}[H]
    \centering
    \includegraphics[width=\linewidth,height=0.68\textheight,keepaspectratio]{#1.png}
    \caption{#2}
    \label{#3}
  \end{figure}
}

% \hinhui{tên_file_không_đuôi}{tiêu đề màn hình}{nhãn}
% Tự động phát hiện ảnh chụp trong thư mục Images/GiaoDien/, nếu chưa có ảnh sẽ hiển thị hộp placeholder nhắc nhở tên file.
\newcommand{\hinhui}[3]{%
  \begin{figure}[H]
    \centering
    \IfFileExists{Images/GiaoDien/#1.png}{%
      \includegraphics[width=0.92\linewidth,keepaspectratio]{Images/GiaoDien/#1.png}%
    }{%
      \begin{tikzpicture}
        \node[draw=blue!60, dashed, fill=blue!5, rounded corners=4pt,
              minimum height=4.2cm, text width=0.88\linewidth, align=center]
          {\textbf{\color{blue!80!black}ẢNH GIAO DIỆN: #1.png}\\[0.25cm]
           \small\color{gray!80}Vui lòng chụp màn hình và lưu vào: \texttt{report/Images/GiaoDien/#1.png}\\[0.15cm]
           \textit{#2}};
      \end{tikzpicture}%
    }
    \caption{#2}
    \label{#3}
  \end{figure}
}

% \hinhdo{ChuongX/ten_file_khong_duoi}{chú thích}{nhãn}
\newcommand{\hinhdo}[3]{%
  \begin{figure}[H]
    \centering
    \includegraphics[width=\linewidth,height=0.82\textheight,keepaspectratio]{#1.png}
    \caption{#2}
    \label{#3}
  \end{figure}%
}

% \hinhdongang{ChuongX/ten_file_khong_duoi}{chú thích}{nhãn} (hiển thị xuôi chiều trang, không xoay 90 độ)
\newcommand{\hinhdongang}[3]{%
  \begin{figure}[H]
    \centering
    \includegraphics[width=\linewidth,height=0.68\textheight,keepaspectratio]{#1.png}
    \caption{#2}
    \label{#3}
  \end{figure}%
}

% \hinherd{duong_dan_file}{chú thích}{nhãn} (dành cho các sơ đồ ERD Chen)
\newcommand{\hinherd}[3]{%
  \begin{figure}[H]
    \centering
    \includegraphics[width=0.96\linewidth,height=0.76\textheight,keepaspectratio]{#1}
    \caption{#2}
    \label{#3}
  \end{figure}%
}
